\SixthArticle{35}
Nicht nur das Produkt $\Pi\lambda.\Psi\lambda$, sondern auch die Funktion $\Psi\lambda$ selbst lässt sich durch ein bestimmtes Integral darstellen. Bezeichnet $k$ eine positive ganze Zahl, so hat das Integral $\int\frac{x^\lambda-x^{\lambda+k}}{1-x}.\mathrm{d}x$ von $x=0$ bis $x=1$ den Wert
\[
=\frac{1}{\lambda+1}+\frac{1}{\lambda+2}+\frac{1}{\lambda+3}+\text{usw.}+\frac{1}{\lambda+k}
\]
Da ferner der Wert des Integrals $\int\left(\frac{1}{1-x}-\frac{kx^{k-1}}{1-x^k}\right)\mathrm{d}x$ allgemein $=\mathrm{Const.}+\log\frac{1-x^k}{1-x}$ ist, ist sein Wert zwischen den Grenzen $x=0$ und $x=1$ $=\log k$; daraus ergibt sich für das Integral $S=\int\left(\frac{1-x^\lambda+x^{\lambda+k}}{1-x}-\frac{kx^{k-1}}{1-x^k}\right)\mathrm{d}x$ zwischen denselben Grenzen der Wert
\[
=\log k-\frac{1}{\lambda+1}-\frac{1}{\lambda+2}-\frac{1}{\lambda+3}-\text{usw.}-\frac{1}{\lambda+k}
\]
Diesen Ausdruck wollen wir mit $\Omega$ bezeichnen. Zerlegen wir das Integral $S$ in zwei Teile:
\[
\int\left(\frac{1-x^\lambda}{1-x}\right)\mathrm{d}x+\int\left(\frac{x^{\lambda+k}}{1-x}-\frac{kx^{k-1}}{1-x^k}\right)\mathrm{d}x
\]
Der erste Teil $\int\frac{1-x^\lambda}{1-x}.\mathrm{d}x$ geht, wenn $x=y^k$ gesetzt wird, über in
\[
\int\frac{ky^{k-1}-ky^{\lambda k+k-1}}{1-y^k}\mathrm{d}y
\]
woraus sich von selbst ergibt, dass sein Wert von $x=0$ bis $x=1$ dem Wert des Integrals
\[
\int\frac{kx^{k-1}-kx^{\lambda k+k-1}}{1-x^k}\mathrm{d}x
\]
zwischen denselben Grenzen gleich ist, da man den Buchstaben $y$ unter dieser Einschränkung offenbar in $x$ ändern darf. Hierdurch wird das Integral $S$ zwischen denselben Grenzen zu
\[
=\int\left(\frac{x^{\lambda+k}}{1-x}-\frac{kx^{\lambda k+k-1}}{1-x^k}\right)\mathrm{d}x
\]
Dieses Integral geht aber, wenn $x^k=z$ gesetzt wird, über in
\SourceLeaf{168}
\[
\int\left(\frac{z^{\frac{\lambda+1}{k}}}{k(1-z)^{\frac{1}{k}}}-\frac{z^\lambda}{1-z}\right)\mathrm{d}z
\]
was also zwischen den Grenzen $z=0$ und $z=1$ genommen gleich $\Omega$ ist. Wächst aber $k$ ins Unendliche, so ist der Grenzwert von $\Omega$ gleich $\Psi\lambda$, der Grenzwert von $\frac{\lambda+1}{k}$ gleich $0$, der Grenzwert von $k(1-z)^{\frac{1}{k}}$ dagegen gleich $\log\frac{1}{z}$ oder $-\log z$. Wir haben daher
\[
\begin{aligned}
&[77]\qquad \Psi\lambda\\[4pt]
&=\int\left(\frac{1}{\log\frac{1}{z}}-\frac{z^\lambda}{1-z}\right)\mathrm{d}z
=\int\left(-\frac{1}{\log z}-\frac{z^\lambda}{1-z}\right)\mathrm{d}z
\end{aligned}
\]
von $z=0$ bis $z=1$.

\SixthArticle{36}
Die bestimmten Integrale, durch die oben die Funktionen $\Pi\lambda$, $\Pi\lambda.\Psi\lambda$ ausgedrückt wurden, mussten auf solche Werte von $\lambda$ beschränkt werden, für die $\lambda+1$ eine positive Größe ist. Diese Einschränkung ergab sich aus der Herleitung selbst; tatsächlich lässt sich leicht erkennen, dass jene Integrale für andere Werte von $\lambda$ stets unendlich werden, obwohl die Funktionen $\Pi\lambda$, $\Pi\lambda.\Psi\lambda$ endlich bleiben können. Auch der Gültigkeit der Formel 77 muss gewiss dieselbe Bedingung zugrunde liegen, dass $\lambda+1$ eine positive Größe ist (andernfalls wird das Integral nämlich gewiss unendlich, obwohl die Funktion $\Psi\lambda$ endlich bleibt); doch scheint die Herleitung der Formel auf den ersten Blick allgemein und keiner Einschränkung unterworfen zu sein. Bei genauerer Betrachtung zeigt sich aber leicht, dass diese Einschränkung bereits in der Analyse selbst enthalten ist, durch die die Formel gewonnen wurde. Wir haben nämlich stillschweigend vorausgesetzt, dass das Integral $\int\frac{1-x^\lambda}{1-x}\mathrm{d}x$, für das wir das gleiche Integral $\int\frac{kx^{k-1}-kx^{\lambda k+k-1}}{1-x^k}\mathrm{d}x$ eingesetzt haben, einen \emph{endlichen} Wert hat; diese Bedingung verlangt, dass $\lambda+1$ eine positive Größe ist. Aus unserer Analyse folgt zwar, dass diese beiden Integrale stets gleich sind, wenn das letztgenannte von $x=0$ bis $x=1-\omega$, das erstgenannte von $x=0$ bis $x=(1-\omega)^k$ erstreckt wird, wie klein die Größe $\omega$ auch sein mag, sofern sie nicht $=0$ ist. Trotzdem konvergieren in dem Fall, in dem $\lambda+1$ keine positive Größe ist, die beiden Integrale von $x=0$ bis zu \emph{derselben} Grenze $x=1-\omega$ keineswegs zur Gleichheit; vielmehr wächst ihre Differenz dann ins Unendliche, wenn $\omega$ gegen Null abnimmt. Dieses Beispiel zeigt, wie viel Umsicht bei der Behandlung unendlicher Größen nötig ist, die nach unserem Urteil in analytischen Schlüssen nur insoweit zugelassen werden dürfen, als sie sich auf die Theorie der Grenzwerte zurückführen lassen.
\SourceLeaf{169}
\SixthArticle{37}
Setzt man in Formel 77 $z=e^{-u}$, so ist ersichtlich, dass sie sich auch folgendermaßen darstellen lässt:
\[
\begin{aligned}
&\Psi\lambda=-\int\left(\frac{e^{-u}}{u}-\frac{e^{-u\lambda-u}}{1-e^{-u}}\right)\mathrm{d}u,
\\[4pt]
&\quad\text{von }u=\infty\text{ bis }u=0,\text{ d. h.}
\end{aligned}
\]
\[
\begin{aligned}
&\text{[78]}\qquad
\Psi\lambda=\int\left(\frac{e^{-u}}{u}-\frac{e^{-\lambda u}}{e^u-1}\right)\mathrm{d}u,
\\[4pt]
&\quad\text{von }u=0\text{ bis }u=\infty.
\end{aligned}
\]
(Ebenso geht der in Artikel 28 angegebene Wert von $\Pi\lambda$, wenn $e^{-y}=\upsilon$ gesetzt wird, in den folgenden über:
\[
\Pi\lambda=\int\left(\log\frac{1}{\upsilon}\right)^\lambda\mathrm{d}\upsilon,
\quad\text{von }\upsilon=0\text{ bis }\upsilon=1)
\]
Ferner ergibt sich aus Formel 77
\[
\text{[79]}\qquad
\Psi\lambda-\Psi\mu=\int\frac{z^\mu-z^\lambda}{1-z}\mathrm{d}z,
\quad\text{von }z=0\text{ bis }z=1
\]
wobei außer $\lambda+1$ auch $\mu+1$ eine positive Größe sein muss.

Setzt man in derselben Formel 77 $z=u^\alpha$, wobei $\alpha$ eine positive Größe bezeichnet, so ergibt sich
\[
\begin{aligned}
&\Psi\lambda=\int\left(-\frac{u^{\alpha-1}}{\log u}
-\frac{\alpha u^{\alpha\lambda+\alpha-1}}{1-u^\alpha}\right)\mathrm{d}u,
\\[4pt]
&\quad\text{von }u=0\text{ bis }u=1
\end{aligned}
\]
und da ebenso für einen positiven Wert von $\beta$
\[
\Psi\lambda=\int\left(-\frac{u^{\beta-1}}{\log u}
-\frac{\beta u^{\beta\lambda+\beta-1}}{1-u^\beta}\right)\mathrm{d}u
\]
gesetzt werden kann, ist ersichtlich, dass
\[
0=\int\left(\frac{u^{\alpha-1}-u^{\beta-1}}{\log u}
+\frac{\alpha u^{\alpha\lambda+\alpha-1}}{1-u^\alpha}
-\frac{\beta u^{\beta\lambda+\beta-1}}{1-u^\beta}\right)\mathrm{d}u
\]
oder
\[
\int\frac{u^{\alpha-1}-u^{\beta-1}}{\log u}\mathrm{d}u
=\int\left(\frac{\beta u^{\beta\lambda+\beta-1}}{1-u^\beta}
-\frac{\alpha u^{\alpha\lambda+\alpha-1}}{1-u^\alpha}\right)\mathrm{d}u
\]
wobei die Integrale stets von $u=0$ bis $u=1$ erstreckt werden. Setzt man aber $\lambda=0$, so kann das letztere Integral \emph{unbestimmt} angegeben werden; es ist nämlich $=\log\frac{1-u^\alpha}{1-u^\beta}$, wenn es für $u=0$ verschwinden soll. Da man also für $u=1$ den Wert $\frac{1-u^\alpha}{1-u^\beta}=\frac{\alpha}{\beta}$ ansetzen muss, ist das Integral $\log\frac{\alpha}{\beta}=\int\frac{u^{\alpha-1}-u^{\beta-1}}{\log u}\mathrm{d}u$ von $u=0$ bis $u=1$, ein Satz, den der berühmte \textsc{Euler} einst mit anderen Methoden hergeleitet hat.
\par\medskip\noindent\hfill\rule{.25\linewidth}{.4pt}\hfill\par
